\begin{abstract}
\noindent 多收集病例能让再生数估计精确多少，取决于要估计什么。本文在负二项传播与二项报告的框架下区分两个估计目标。估计当期再生数时，报告数除以预计报告规模得到的估计量在每个有限规模上都达到 Cramér--Rao 下界，方差随规模趋于零。估计跨期平均传播水平时，方差为 $v_R+A/m$，其中 $v_R$ 是环境波动造成的方差，$m$ 是预计报告病例数；环境项构成规模无法突破的底板。两项相等处的交叉点 $m_\times=A/v_R$ 为规模提供了一把尺子。规模达到 $m_\times$ 后，再翻一倍只能把方差降低 25\%。这一结构有三点设计含义。扩大覆盖面比提高报告率更能压低方差。多期预算下，最优单期规模取成本比与交叉点的几何平均。小规模时应使用精确增长检验，常用的正态近似检验会超过名义水平。更新过程仿真表明，在时间序列中单代公式只能作参照曲线，不构成误差下界。本文用美国国家医疗安全网络 16,218 个州周的流感住院数据做检验。流行初期的预测误差呈“常数加规模反比”的形状，经验交叉点为 22.6 例/周，与理论情景区间 21.3--49.2 同一量级。规模项在多种替代规格和两组外部数据中都为正。拟合的截距与斜率则都在情景隐含值的三倍以上。滞后权重的比较表明，报告噪声使分母随机，抬高了斜率；截距随分母的回溯滞后上升，与环境的时间相关一致。因此经验交叉点只能作量级参考。达峰期 88.9\% 的误差来自基准偏移，换用阶段自身基准后规模效应依然存在。结果提示，预警和趋势估计宜分开设计。

\vspace{0.6em}
\noindent\textbf{关键词：}再生数；环境随机性；超传播；监测规模；方差分解；流感住院
\end{abstract}
